\chapter{流体运动学}

\intro{流体运动学（fluid kinematics）研究流体运动的几何性质，而不涉及引起运动的力。运动学是连接流体静力学和流体动力学的桥梁。本章将系统地讨论描述流体运动的数学方法、流体微团的运动分解、涡量动力学基础以及流动的分类。}

\section{描述流体运动的两种方法}

\subsection{拉格朗日方法与欧拉方法回顾}

如第一章所述，描述流体运动有两种基本方法。拉格朗日方法以流体质点为研究对象，追踪其位置、速度等物理量随时间的变化。设流体质点在初始时刻 \(t_0\) 的坐标为 \(\mathbf{X}\)（物质坐标），则质点在任意时刻 \(t\) 的位置为：

\begin{myformula}
\mathbf{x} = \mathbf{x}(\mathbf{X}, t)
\end{myformula}

质点的速度和加速度分别为：
\begin{myformula}
\mathbf{u} = \left(\frac{\partial \mathbf{x}}{\partial t}\right)_{\mathbf{X}},\quad
\mathbf{a} = \left(\frac{\partial^2 \mathbf{x}}{\partial t^2}\right)_{\mathbf{X}}
\end{myformula}

欧拉方法以空间固定点为研究对象，考察通过空间各点的流体质点的物理量随时间的变化。速度场 \(\mathbf{u}(\mathbf{x}, t)\) 是欧拉描述中的核心物理量。

\begin{myformula}
\mathbf{u} = \mathbf{u}(\mathbf{x}, t) = (u_x(x,y,z,t), u_y(x,y,z,t), u_z(x,y,z,t))
\end{myformula}

加速度由物质导数给出：
\begin{myformula}
\mathbf{a} = \frac{D\mathbf{u}}{Dt} = \frac{\partial \mathbf{u}}{\partial t} + \mathbf{u}\cdot\Nabla\mathbf{u}
\end{myformula}

两种描述之间可以通过速度场建立联系。若已知欧拉速度场 \(\mathbf{u}(\mathbf{x}, t)\)，求解微分方程组 \(d\mathbf{x}/dt = \mathbf{u}(\mathbf{x}, t)\) 并代入初始条件 \(\mathbf{x}(t_0) = \mathbf{X}\)，即可得到拉格朗日描述 \(\mathbf{x} = \mathbf{x}(\mathbf{X}, t)\)。

\subsection{流线、迹线与流管}

流体运动中有三种重要的曲线用于可视化流场结构。

迹线（pathline）：同一流体质点在不同时刻所经过的位置连成的曲线。迹线是拉格朗日描述下的概念，满足：
\begin{myformula}
\frac{d\mathbf{x}}{dt} = \mathbf{u}(\mathbf{x}, t),\quad \mathbf{x}(t_0) = \mathbf{X}
\end{myformula}

流线（streamline）：某一瞬时处处与速度矢量相切的曲线。流线是欧拉描述下的概念，满足：
\begin{myformula}
\frac{dx}{u_x} = \frac{dy}{u_y} = \frac{dz}{u_z}
\end{myformula}

流线具有以下性质：
\begin{enumerate}
    \item 除速度为零的点（驻点，stagnation point）外，流线不能相交。
    \item 流线的疏密反映速度的大小：流线密集处流速大，稀疏处流速小。
    \item 在定常流动中，流线与迹线重合。
\end{enumerate}

\begin{figure}[H]
\centering
\begin{tikzpicture}[scale=1.3]
    \draw[->] (0,0) -- (5,0) node[right] {$x$};
    \draw[->] (0,0) -- (0,4) node[above] {$y$};
    \foreach \y in {0.5,1.0,...,3.5} {
        \draw[blue,->,domain=0:4.5,smooth,samples=30]
            plot(\x, {\y + 0.5*sin(1.5*\x r)});
    }
    \draw[red,thick,->,domain=0:4,smooth,samples=50]
        plot(\x, {2.0 + 0.8*sin(1.5*\x r)});
    \node[red] at (4.5,2.8) {迹线};
    \node[blue] at (4.5,0.7) {流线族};
    \fill[red] (0,2.0) circle (0.08) node[left] {$\mathbf{X}$};
\end{tikzpicture}
\caption{非定常流动中的流线（蓝色）与迹线（红色）对比}
\label{fig:streamline-pathline}
\end{figure}

流管（stream tube）：在流场中任取一条不与流线重合的封闭曲线 \(C\)，通过 \(C\) 上各点的流线所围成的管状曲面称为流管。由于流体质点的速度始终沿流线方向，流体不能穿过流管壁面。因此，流管类似于一个无形的固体管道，在定常流动中流体在流管内运动如同在真实的管道中一样。

烟线（streakline）：先后通过空间同一点的所有流体质点连成的线。在实验流场可视化中，通过固定点连续注入染料或烟雾得到的就是烟线。在定常流动中，流线、迹线和烟线三者重合。

\section{速度场的分解}

\subsection{亥姆霍兹速度分解定理}

流体微团在运动过程中不仅要平动，还可能发生旋转和变形。亥姆霍兹（Helmholtz）速度分解定理揭示了流体微团运动的完整图形。

考虑流场中邻近两点 \(M(\mathbf{x})\) 和 \(M'(\mathbf{x} + d\mathbf{x})\)，其速度分别为 \(\mathbf{u}\) 和 \(\mathbf{u} + d\mathbf{u}\)。由多元函数的全微分：

\begin{myformula}
du_i = \frac{\partial u_i}{\partial x_j} dx_j
\end{myformula}

速度梯度张量 \(\partial u_i/\partial x_j\) 可以分解为对称部分和反对称部分：

\begin{myformula}
\frac{\partial u_i}{\partial x_j} = \frac{1}{2}\left(\frac{\partial u_i}{\partial x_j} + \frac{\partial u_j}{\partial x_i}\right) + \frac{1}{2}\left(\frac{\partial u_i}{\partial x_j} - \frac{\partial u_j}{\partial x_i}\right) \equiv e_{ij} + \Omega_{ij}
\end{myformula}

其中：
\[
e_{ij} = \frac{1}{2}\left(\frac{\partial u_i}{\partial x_j} + \frac{\partial u_j}{\partial x_i}\right)
\]
为应变率张量（rate-of-strain tensor），是一个对称张量（\(e_{ij}=e_{ji}\)），描述流体微团的变形运动。

\[
\Omega_{ij} = \frac{1}{2}\left(\frac{\partial u_i}{\partial x_j} - \frac{\partial u_j}{\partial x_i}\right)
\]
为旋转率张量（rotation rate tensor），是一个反对称张量（\(\Omega_{ij}=-\Omega_{ji}\)），描述流体微团的旋转运动。

\subsection{亥姆霍兹定理的物理意义}

因此，邻近点 \(M'\) 相对于 \(M\) 的速度可以写为：

\begin{myTheorem}{亥姆霍兹速度分解定理}
流体微团的运动可以分解为三部分：
\begin{myformula}
\mathbf{u}(\mathbf{x}+d\mathbf{x}) = \mathbf{u}(\mathbf{x}) + \boldsymbol{\Omega} \times d\mathbf{x} + \mathbf{e} \cdot d\mathbf{x}
\end{myformula}
\begin{enumerate}
    \item 平动（translation）：以中心速度 \(\mathbf{u}(\mathbf{x})\) 整体平移。
    \item 旋转（rotation）：以角速度 \(\boldsymbol{\Omega} = \frac{1}{2}\Nabla\times\mathbf{u}\) 绕某瞬时轴旋转（刚体转动）。
    \item 变形（deformation）：由应变率张量 \(e_{ij}\) 描述的纯变形（线变形和角变形）。
\end{enumerate}
\end{myTheorem}

反对称张量 \(\Omega_{ij}\) 与涡量矢量 \(\boldsymbol{\omega} = \Nabla\times\mathbf{u}\) 有如下关系：
\begin{myformula}
\Omega_{ij} = -\frac{1}{2}\varepsilon_{ijk}\omega_k
\end{myformula}

即涡量 \(\boldsymbol{\omega}\) 是旋转角速度的两倍。

\begin{figure}[H]
\centering
\begin{tikzpicture}[scale=1.5]
    \draw[->] (-1,0) -- (5,0) node[right] {$x$};
    \draw[->] (0,-1) -- (0,4) node[above] {$y$};
    \draw[thick] (1,2) rectangle (3,3);
    \node at (2,2.5) {$M$};
    \draw[->] (2,2.5) -- (2.8,2.5) node[above] {平动};
    \begin{scope}[shift={(4,0)}]
        \draw[thick] (1,2) rectangle (3,3);
        \draw[thick,dashed,rotate around={15:(2,2.5)}] (1,2) rectangle (3,3);
        \node at (2,2.5) {旋转};
    \end{scope}
    \begin{scope}[shift={(0,-2)}]
        \draw[thick] (1,2) rectangle (3,3);
        \draw[thick,dashed] (1.3,1.8) rectangle (3.2,3.2);
        \node at (2.3,2.5) {变形};
    \end{scope}
    \begin{scope}[shift={(4,-2)}]
        \draw[thick] (1,2) rectangle (3,3);
        \draw[thick,dashed] (1.5,2) -- (2.5,3);
        \node at (2,2.5) {剪切};
    \end{scope}
\end{tikzpicture}
\caption{流体微团运动的分解：平动、旋转、线变形与剪切变形}
\label{fig:helmholtz}
\end{figure}

\section{应变率张量}

\subsection{线应变率与角应变率}

应变率张量的对角线分量表示线应变率（elongation rate）：

\begin{myformula}
e_{11} = \frac{\partial u_1}{\partial x_1},\quad
e_{22} = \frac{\partial u_2}{\partial x_2},\quad
e_{33} = \frac{\partial u_3}{\partial x_3}
\end{myformula}

对于不可压缩流动，由连续性方程 \(\partial u_i/\partial x_i = 0\)，故应变率张量的迹恒为零：
\begin{myformula}
e_{ii} = \frac{\partial u_1}{\partial x_1} + \frac{\partial u_2}{\partial x_2} + \frac{\partial u_3}{\partial x_3} = 0
\end{myformula}

应变率张量的非对角线分量表示角应变率（shear rate）：

\begin{myformula}
e_{12} = e_{21} = \frac{1}{2}\left(\frac{\partial u_1}{\partial x_2} + \frac{\partial u_2}{\partial x_1}\right)
\end{myformula}

\begin{myformula}
e_{13} = e_{31} = \frac{1}{2}\left(\frac{\partial u_1}{\partial x_3} + \frac{\partial u_3}{\partial x_1}\right)
\end{myformula}

\begin{myformula}
e_{23} = e_{32} = \frac{1}{2}\left(\frac{\partial u_2}{\partial x_3} + \frac{\partial u_3}{\partial x_2}\right)
\end{myformula}

\subsection{应变率张量的主值}

作为实对称张量，应变率张量可以通过施密特正交变换化为对角型。其三个主值 \(e^{(1)}\)、\(e^{(2)}\)、\(e^{(3)}\) 满足特征方程：
\begin{myformula}
\det(e_{ij} - e\delta_{ij}) = 0
\end{myformula}

主方向上的三个互相垂直的线元只有线应变率而无角应变率。对于不可压缩流动，三个主值之和为零。

应变率张量在牛顿流体的本构关系中起到核心作用。牛顿流体的偏应力张量与应变率张量成正比：

\begin{myformula}
\tau_{ij} = 2\mu e_{ij}
\end{myformula}

这是纳维—斯托克斯方程中粘性项的来源。

\subsection{涡量与速度梯度张量的不变量}

速度梯度张量 \(\partial u_i/\partial x_j\) 有三个主不变量。第一不变量（散度）：
\[
P = \frac{\partial u_i}{\partial x_i} = \Nabla\cdot\mathbf{u}
\]

第二不变量常用于湍流的涡辨识（如 \(Q\)-criterion）：
\begin{myformula}
Q = \frac{1}{2}(\Omega_{ij}\Omega_{ij} - e_{ij}e_{ij})
\end{myformula}

在 \(Q > 0\) 的区域，旋转占主导地位，表明存在涡结构。

\section{涡量}

\subsection{涡量的定义与物理意义}

涡量（vorticity）定义为速度场的旋度：

\begin{mydefinition}{涡量}
\begin{myformula}
\boldsymbol{\omega} = \Nabla \times \mathbf{u}
\end{myformula}
在直角坐标系中：
\begin{myformula}
\boldsymbol{\omega} = \left(\frac{\partial u_z}{\partial y} - \frac{\partial u_y}{\partial z},\;
\frac{\partial u_x}{\partial z} - \frac{\partial u_z}{\partial x},\;
\frac{\partial u_y}{\partial x} - \frac{\partial u_x}{\partial y}\right)
\end{myformula}
\end{mydefinition}

涡量的物理意义是流体微团旋转角速度的两倍。流体微团绕自身形心的平均旋转角速度为 \(\boldsymbol{\Omega} = \boldsymbol{\omega}/2\)。

涡量场的基本性质：
\begin{enumerate}
    \item 涡量场是无源场（管式场）：\(\Nabla\cdot\boldsymbol{\omega} = \Nabla\cdot(\Nabla\times\mathbf{u}) = 0\)。
    \item 涡量场的无源性意味着涡线不能起止于流体内部，只能形成闭合曲线或终止于边界。
\end{enumerate}

\subsection{涡线、涡管与涡丝}

涡线（vortex line）：某一瞬时处处与涡量矢量相切的曲线。涡线的微分方程为：
\begin{myformula}
\frac{dx}{\omega_x} = \frac{dy}{\omega_y} = \frac{dz}{\omega_z}
\end{myformula}

涡管（vortex tube）：通过一条非涡线的封闭曲线上各点的涡线所围成的管状区域。

涡丝（vortex filament）：截面面积趋近于零的涡管。

通过涡管截面的涡通量（vortex flux）定义为：
\begin{myformula}
\Gamma_{\text{vortex}} = \iint_S \boldsymbol{\omega} \cdot \mathbf{n} \, dS
\end{myformula}

由涡量场的无源性可以证明，涡管强度沿管长不变（亥姆霍兹涡管定理）：同一时刻，通过涡管各截面的涡通量相等。这意味着涡管不能在流体中中断——它要么形成闭合环，要么终止于自由表面或固体壁面。

\subsection{环量与斯托克斯定理}

速度环量（circulation）定义为速度沿封闭曲线的线积分：

\begin{mydefinition}{速度环量}
\begin{myformula}
\Gamma = \oint_C \mathbf{u} \cdot d\mathbf{r}
\end{myformula}
\end{mydefinition}

环量的物理意义是封闭曲线所围区域内涡旋强度的度量。根据斯托克斯定理，速度环量等于通过以该曲线为边界的任意曲面的涡通量：

\begin{myformula}
\Gamma = \oint_C \mathbf{u} \cdot d\mathbf{r} = \iint_S (\Nabla\times\mathbf{u}) \cdot \mathbf{n} \, dS = \iint_S \boldsymbol{\omega} \cdot \mathbf{n} \, dS
\end{myformula}

这一关系是涡动力学中最为重要的基本关系之一，将速度场的整体性质（环量）与局部性质（涡量）联系起来。

\section{有旋流动与无旋流动}

\subsection{无旋流动与速度势}

如果流场中处处 \(\boldsymbol{\omega} = \Nabla\times\mathbf{u} = \mathbf{0}\)，则称该流动为无旋流动（irrotational flow）或势流（potential flow）。

对于无旋流动，速度场可以表示为某一标量函数 \(\phi(\mathbf{x},t)\) 的梯度：

\begin{myTheorem}{速度势}
在无旋流动中，存在标量函数 \(\phi(\mathbf{x},t)\)，使得：
\begin{myformula}
\mathbf{u} = \Nabla\phi
\end{myformula}
\(\phi\) 称为速度势（velocity potential）。
\end{myTheorem}

这是因为旋度为零的充分必要条件是该矢量场可表示为标量函数的梯度（在单连通区域中）。

对于不可压缩的无旋流动，将 \(\mathbf{u} = \Nabla\phi\) 代入连续性方程 \(\Nabla\cdot\mathbf{u} = 0\)，得到：

\begin{myformula}
\Nabla^2\phi = \frac{\partial^2\phi}{\partial x^2} + \frac{\partial^2\phi}{\partial y^2} + \frac{\partial^2\phi}{\partial z^2} = 0
\end{myformula}

即速度势满足拉普拉斯方程。这表明不可压缩无旋流动的求解归结为求解拉普拉斯方程（椭圆型偏微分方程），这一问题具有良好的数学性质（解的存在性、唯一性、叠加原理适用等）。

\begin{remark}
无旋假设在流体力学中是一个非常有用的近似。在大雷诺数流动中，粘性效应局限于薄边界层和尾流中，主流区的流动可以近似为无旋流动。这使得我们可以利用势流理论有效地解决许多实际工程问题。
\end{remark}

\subsection{有旋流动的特征}

当 \(\boldsymbol{\omega} \neq \mathbf{0}\) 时，流动为有旋流动（rotational flow）。有旋流动具有以下特征：

\begin{itemize}
    \item 不能定义单值的速度势函数。
    \item 流体微团具有旋转角速度。
    \item 涡量与速度场之间存在复杂的非线性相互作用（涡拉伸、涡倾斜等机制）。
    \item 即使初始无旋的流动，由于粘性或斜压效应也可能产生涡量。
\end{itemize}

在自然和工程中，有旋流动非常普遍。边界层中的流动、分离流动、漩涡脱落（如卡门涡街）、湍流等都是典型的有旋流动。

\begin{figure}[H]
\centering
\begin{tikzpicture}[scale=1.2]
    \begin{scope}[local bounding box=irrot]
        \draw[->] (0,0) -- (4,0) node[right] {$x$};
        \draw[->] (0,0) -- (0,3) node[above] {$y$};
        \foreach \y in {0.8,1.6,2.4} {
            \draw[blue,->,domain=0:3.5,smooth,samples=20]
                plot(\x, \y);
        }
        \draw[red,thick] (1.5,1.6) -- (2.0,1.6) -- (2.0,2.0) -- (1.5,2.0) -- cycle;
        \node at (2,-0.5) {无旋流动（平行剪切流）};
    \end{scope}
    \begin{scope}[shift={(5,0)},local bounding box=rot]
        \draw[->] (0,0) -- (4,0) node[right] {$x$};
        \draw[->] (0,0) -- (0,3) node[above] {$y$};
        \foreach \r in {0.3,0.7,1.1,1.5} {
            \draw[blue] (0,0) circle (\r);
        }
        \draw[->,red,thick] (1.5,1.5) arc[start angle=45, end angle=135, radius=0.5];
        \node at (2,-0.5) {有旋流动（强制涡）};
    \end{scope}
\end{tikzpicture}
\caption{无旋流动与有旋流动对比}
\label{fig:irrot-vs-rot}
\end{figure}

\section{流函数}

\subsection{平面流动的流函数}

对于二维不可压缩平面流动，速度分量 \(u = u(x,y)\)、\(v = v(x,y)\) 满足连续性方程：
\[
\frac{\partial u}{\partial x} + \frac{\partial v}{\partial y} = 0
\]

根据这个条件，可以引入一个标量函数 \(\psi(x,y)\)（流函数，stream function），使得速度分量自动满足连续性方程：

\begin{mydefinition}{平面流动的流函数}
对于二维不可压缩流动，存在流函数 \(\psi(x,y)\)，满足：
\begin{myformula}
u = \frac{\partial \psi}{\partial y},\quad v = -\frac{\partial \psi}{\partial x}
\end{myformula}
\end{mydefinition}

流函数的物理意义和性质：
\begin{enumerate}
    \item 沿流线 \(\psi = \text{const}\)，等流函数线即为流线。两条流线间的体积流量等于两者流函数值之差。
    \item 涡量的 \(z\) 分量与流函数的关系：\(\omega_z = \frac{\partial v}{\partial x} - \frac{\partial u}{\partial y} = -\Nabla^2\psi\)。
    \item 对于平面无旋流动，流函数也满足拉普拉斯方程 \(\Nabla^2\psi = 0\)。
\end{enumerate}

对于平面无旋流动，速度势 \(\phi\) 和流函数 \(\psi\) 满足柯西—黎曼（Cauchy-Riemann）条件：
\begin{myformula}
\frac{\partial\phi}{\partial x} = \frac{\partial\psi}{\partial y},\quad
\frac{\partial\phi}{\partial y} = -\frac{\partial\psi}{\partial x}
\end{myformula}

这表明 \(\phi\) 和 \(\psi\) 可以构成一个解析函数的实部和虚部，从而可以运用复变函数论的强大工具分析平面势流。等势线（\(\phi=\text{const}\)）和等流函数线（\(\psi=\text{const}\)）构成一组正交曲线网。

\subsection{轴对称流动的流函数}

对于轴对称流动（如在球坐标系或柱坐标系中描述的绕回转体的流动），在连续性方程的约束下，同样可以引入流函数（斯托克斯流函数）。

在柱坐标系 \((r, \theta, z)\) 中，对于轴对称（\(\partial/\partial\theta = 0\)）的不可压缩流动，连续性方程为：
\[
\frac{1}{r}\frac{\partial}{\partial r}(r u_r) + \frac{\partial u_z}{\partial z} = 0
\]

引入斯托克斯流函数 \(\psi(r,z)\)：
\begin{myformula}
u_r = -\frac{1}{r}\frac{\partial\psi}{\partial z},\quad
u_z = \frac{1}{r}\frac{\partial\psi}{\partial r}
\end{myformula}

在球坐标系 \((R, \theta, \varphi)\) 中，对于轴对称不可压缩流动：
\begin{myformula}
u_R = \frac{1}{R^2\sin\theta}\frac{\partial\psi}{\partial\theta},\quad
u_\theta = -\frac{1}{R\sin\theta}\frac{\partial\psi}{\partial R}
\end{myformula}

\subsection{三维流动的流函数}

对于一般三维不可压缩流动，可以引入矢量势（vector potential）\(\boldsymbol{\Psi}\)，使得：
\begin{myformula}
\mathbf{u} = \Nabla\times\boldsymbol{\Psi}
\end{myformula}

这自动满足连续性方程 \(\Nabla\cdot\mathbf{u} = 0\)。但矢量势的选择不是唯一的（规范自由度），通常要求 \(\Nabla\cdot\boldsymbol{\Psi} = 0\)（库伦规范）。这种表示在实际计算中使用较少，主要见于理论分析。

\section{流动的分类}

\subsection{定常与非定常流动}

根据流动参数是否随时间变化，流动分为定常流动（steady flow）和非定常流动（unsteady flow）。

在定常流动中，\(\partial/\partial t = 0\)（对任何物理量），速度场仅为空间坐标的函数 \(\mathbf{u} = \mathbf{u}(\mathbf{x})\)。在定常流动中，流线、迹线和烟线重合，流动的几何图像不随时间改变。

非定常流动中至少有一部分物理量随时间变化。自然界中的绝大多数流动都是非定常的，但在许多工程问题中，当时间平均量的变化足够缓慢时，可采用准定常假设（quasi-steady assumption）进行简化。

非定常流动的例子包括：
\begin{itemize}
    \item 启动/停止过程中的流动
    \item 周期性流动（如圆柱后缘的卡门涡街）
    \item 活塞式发动机中的脉动流动
    \item 波浪运动
    \item 湍流的瞬时运动
\end{itemize}

\subsection{均匀与非均匀流动}

均匀流动（uniform flow）是指速度矢量不随空间位置变化的流动（但可随时间变化），即 \(\Nabla\mathbf{u} = 0\)。均匀流动中所有流体质点作相同的平动，没有变形和旋转。

非均匀流动（non-uniform flow）中速度场随空间位置变化。实际流动绝大多数都是非均匀的。如果速度沿某一方向的变化率远小于沿其他方向的变化率，则可以对该方向作均匀近似。

\subsection{一维、二维与三维流动}

根据速度场对空间坐标的依赖程度，流动可分为：

一维流动：速度只依赖于一个空间坐标，如充分发展的管流（泊肃叶流动）\(u = u(r)\)。

二维流动：速度只依赖于两个空间坐标，如平面流动 \(u=u(x,y), v=v(x,y), w=0\)，或轴对称流动 \(u_r=u_r(r,z), u_z=u_z(r,z), u_\theta=0\)。

三维流动：速度依赖于全部三个空间坐标。真实流动本质上都是三维的，但根据几何和流动条件的对称性可简化为二维或一维问题进行分析。

\subsection{层流与湍流概述}

根据流动的形态，可分为层流（laminar flow）和湍流（turbulent flow）。

层流：流体分层流动，各层之间互不掺混，流动有序而规则。层流中粘性力起主导作用，流动对扰动的阻尼大。

湍流：流体微团作不规则、随机的脉动运动，各层之间剧烈掺混。湍流中惯性力起主导作用，流动对扰动敏感。

流动状态由雷诺数（Reynolds number）判定：

\begin{myformula}
Re = \frac{U L}{\nu}
\end{myformula}

其中 \(U\) 为特征速度，\(L\) 为特征长度，\(\nu\) 为运动粘度。

对于圆管流动，当 \(Re < 2300\) 时为层流，\(Re > 4000\) 时为湍流，中间为过渡区。对于外部绕流（如平板边界层），转捩雷诺数约为 \(Re \approx 5\times 10^5\)。

湍流的详细研究属于高等流体力学的范畴，本课程在基础部分主要讨论层流和势流。

\subsection{可压缩与不可压缩流动}

根据密度是否变化，流动分为可压缩流动和不可压缩流动。

不可压缩流动：\(\rho = \text{const}\)（或 \(D\rho/Dt = 0\)），连续性方程简化为 \(\Nabla\cdot\mathbf{u} = 0\)。液体在多数情况下可视为不可压缩，低速气体流动也常作不可压缩假设。

可压缩流动：密度随压强和温度变化。当马赫数 \(Ma = U/c > 0.3\) 时（\(c\) 为声速），压缩性效应不可忽略。可压缩流动中可能出现激波、膨胀波等独特现象。

\subsection{本章小结}

本章系统地阐述了流体运动学的基本概念和理论。我们回顾了拉格朗日描述和欧拉描述的对比，明确了流线、迹线和流管的定义与性质。亥姆霍兹速度分解定理揭示了流体微团运动的完整图像——平动、旋转和变形的叠加。应变率张量和涡量是描述流体微团变形和旋转的核心物理量。涡线、涡管的概念以及环量与涡通量之间的关系（斯托克斯定理）为涡动力学奠定了基础。无旋流动中引入的速度势使问题简化为求解拉普拉斯方程。流函数的概念（特别是平面流动中与复变函数理论的联系）提供了强大的分析工具。最后，我们按不同标准对流动进行了分类，建立了流动类型的系统框架。

\begin{myalgorithm}{二维流场可视化：流线与涡量计算}
\begin{mycode}
import numpy as np
import matplotlib.pyplot as plt

def stream_function(x, y, U0=1.0, Gamma=2.0):
    """点涡 + 均匀流的流函数"""
    r = np.sqrt(x**2 + y**2)
    r = np.maximum(r, 1e-10)
    psi = U0 * y - Gamma / (2*np.pi) * np.log(r)
    return psi

def vorticity(x, y, dx=1e-6):
    """数值计算二维涡量 omega_z = dv/dx - du/dy"""
    _, v_p = velocity(x+dx, y)
    _, v_m = velocity(x-dx, y)
    u_t, _ = velocity(x, y+dx)
    u_b, _ = velocity(x, y-dx)
    omega = (v_p - v_m)/(2*dx) - (u_t - u_b)/(2*dx)
    return omega

def velocity(x, y, U0=1.0, Gamma=2.0):
    r2 = x**2 + y**2
    r2 = np.maximum(r2, 1e-10)
    u = U0 - Gamma/(2*np.pi) * y/r2
    v = Gamma/(2*np.pi) * x/r2
    return u, v

x = np.linspace(-4, 4, 100)
y = np.linspace(-3, 3, 80)
X, Y = np.meshgrid(x, y)
psi = stream_function(X, Y)
levels = np.linspace(-5, 5, 30)
plt.contour(X, Y, psi, levels=levels, colors='blue')
plt.axis('equal')
plt.xlabel('$x$')
plt.ylabel('$y$')
plt.title('点涡 + 均匀流的流线')
plt.show()
\end{mycode}
\end{myalgorithm}
